%%% MA_14_Distribution_Welfare_State.tex -- reorganised; frames moved from the old MA_xx files (see comments)

\section{Distribution and the Welfare State}

%%% [moved from MA_11_Social_Policy.tex, frame 0]
% \redsection{Inequality I: Accumulation, Labour and Capital, Taxes and Transfers}
% \redsection{Inequality II: Social Policy, Education, Pensions, Welfare State}


\begin{frame}{Inequality}
    \mytext{
    \begin{itemize}
        \item inequality of factor reimbursement:
        \begin{itemize}
            \item labour (wages)
            \item capital (interest rates)
        \end{itemize}
        \item inequality of personal income:
        \begin{itemize}
            \item life chances, education
            \item unskilled workers
            \item skilled workers
            \item managers
            \item capital gains
        \end{itemize}
    \end{itemize}
    }
\end{frame}

%%% [moved from MA_11_Social_Policy.tex, frame 15]
\begin{frame}{Income Distribution}
    \mycomment{
    Wage or salaries per hour vary considerably. The top 10\% earners can earn  30\% of all labour income in an economy - check real world data.
    \begin{itemize}
        \item higher tax burden is not redistribution
        \item progressive tax on wages distorts labour markets
        \item fiscal redistribution (payout to lower earners) leaves labour markets unchanged
        \item[$\rightarrow$] efficient
    \end{itemize}
    }
    \showmytikz{../shared/tikz/tikz-incomeshare}

\end{frame}

%%% [moved from MA_11_Social_Policy.tex, frame 16]
\begin{frame}{Wealth Distribution}
    \mycomment{
    Wealth distribution (ownership of assets) is harder to measure. Wealth typically produces income via capital gains. Apart from income, wealth also means control of assets, and is hence much closer linked to power than income.

    There are wealth billionaires, but not (or few) income billionaires.

    Inequality in wealth was at a historic low after WWII and has been sharply increasing in the last decades.
    }
    \showmytikz{../shared/tikz/tikz-wealthdistribution}
\end{frame}

%%% [moved from MA_11_Social_Policy.tex, frame 17]
\begin{frame}{Elephant Curve (Lakner-Milanovic Curve}
    \mycomment{
    In the age of globalisation, incomes of the top earners have increased, while those of the lower-middle class have stagnated.
    }
    \showmytikz{../shared/tikz/tikz-elephantcurve}
\end{frame}

%%% [moved from MA_11_Social_Policy.tex, frame 18]
\begin{frame}{Pareto Distribution}
    \mycomment{
    Vilfredo \noun{Pareto} first observed that both income and wealth distribution follow a power law, and that this is (realtively) stable over time.

    Sometimes the rule \enquote{20\%/80\%} is called \enquote{Pareto Principle}, as roughly 20\% hold 80\% of wealth; but the principle seems to hold for other things as well.

    The figure is illustrative, but realistic; check real world data for your economy.
    }
    \showmytikz{../shared/tikz/tikz-paretodistribution}
\end{frame}

\begin{frame}{Kuznets Curve}
    \showmytikz[1]{../shared/tikz/tikz-kuznets}
    \mycomment{
    In the 50ies/60ies, Simon \noun{Kuznets} argued inequality to rise with income per capita and then fall later. The left half of this curve seems realistic: inequality grew with income per capita. The right half is less plausible. Inequality in OECD countries fell after the Second World War, but is today again on the rise since the 80ies/90ies. 

    The debate was recently resumed by Thomas \noun{Picketty} \cite{piketty2013capital} \cite{piketty2015economie}, arguing that both fall of inequality and rise of income per capita stem from social bargaining rather than laissez-faire economic laws.
    }
\end{frame}

%%% [moved from MA_11_Social_Policy.tex, frame 1]
\begin{frame}{Welfare State}
    \mytext{
    \begin{itemize}
        \item Universal Health Insurance
        \item Child Care
        \item Compulsory Schooling
        \item Public Universities
        \item Unemployment Insurance
        \item Pensions
    \end{itemize}
    }
\end{frame}

%%% [moved from MA_2_Labour_Market.tex, frame 0]

%%% first two moved to Session 9 Labour Market

%%% [moved from MA_2_Labour_Market.tex, frame 2]
\begin{frame}{Employment and Transfers}
    \mycomment{
    Only those employed have an income from labour. However, eveyone was young once and will be old one day; everyone employed now risks unemployment or illness.
    }
    \showmytikz[1,2]{../shared/tikz/tikz-redistribution-realistic}
\end{frame}

%%% [moved from MA_2_Labour_Market.tex, frame 3]
\begin{frame}{Employment and Transfers}
    \mycomment{
    It would be desirable to allow an income to those groups of the population who don't have labour income. 

    This can be achieved by many different means:

    \begin{itemize}
        \item parents are obliged to care for their children. The state assists with money, schooling etc.
        \item unemployment \enquote{insurance} in various forms
        \item health care, universal credit and other means of support for those unable to work
        \item pensions, which can be organised in a variety of ways.
    \end{itemize}
    
    But where can this income come from?

    
    }
    \showmytikz[1,2,3]{../shared/tikz/tikz-redistribution-realistic}
\end{frame}

%%% [moved from MA_2_Labour_Market.tex, frame 4]
\begin{frame}{Employment and Transfers}
    \mycomment{
    Those who have an income must transfer a part of this to those in needs. This can be done via various channels or mechanisms. 
    
    One is insurance, such as unemployment insurance, health insurance or private pension schemes.

    Another is taxation, in which case the labour income is paid to the state, who then provides for those in need.
    In reality the \enquote{split} of contributions between employers and employs is illusory: the factor labour is taxed.

    Another is privately, voluntarily or obligatory, such as parents being obliged to finance their children, or relatives to take care for their family.
    }
    \showmytikz[1,2,3,4]{../shared/tikz/tikz-redistribution-realistic}
\end{frame}

%%% [moved from MA_2_Labour_Market.tex, frame 5]
\begin{frame}{Employment and Transfers}
    \mycomment{
    There is not only labour income, but also capital income, which can be redistributed as well. Capital income is more anonymous, thus would not work as well for direct insurance redistribution.

    Some pension schemes consist of building up capital ownership rather than of direct redistribution.
    }
    \showmytikz[1,2,3,4,5,6]{../shared/tikz/tikz-redistribution-realistic}
\end{frame}

%%% [moved from MA_2_Labour_Market.tex, frame 6]
\begin{frame}{Employment and Transfers}
    \mycomment{
    Not every work done and paid for is officially registered and taxed. Some economies have quite large informal sectors. 
    
    It is by nature difficult to know about the size of wages in this sector. Expect employers and employees to \enquote{share} the taxes \enquote{saved}.
    
    In poor countries, this sector can make up a significant part of the economy.

    Non-official income is by its nature not taxed and does thus not contribute to the formal redistribution schemes.

    There is also an informal capital income sector. There is tax evasion, such as income transferred abroad to tax haves. There can also be undeclared (illicit) income.
    }
    \showmytikz[1,2,3,4,5,6,7,8]{../shared/tikz/tikz-redistribution-realistic}
\end{frame}

%%% [moved from MA_2_Labour_Market.tex, frame 7]
\begin{frame}{Employment and Transfers}
    \mytext{
    \begin{itemize}
        \item The sum of all transfers received must equal the sum of transfers paid (balance mechanic)
        \item otherwise, debt builds up
        \item this is independent of the actual mechanism of redistribution
        \item there are also issues of redistribution within earners of income (high income vs. low income)
        \item the (possible) contribution of capital gains is often overlooked
        \item this only looks at income and consumption (flows), the buildup of wealth (stocks) is out of the picture
    \end{itemize}
    }
\end{frame}

%%% [moved from MA_11_Social_Policy.tex, frame 2]
\begin{frame}{Average vs. Marginal Tax Rates}
    \showmytikz{../shared/tikz/tikz-marginaltaxrates}
\end{frame}

%%% [moved from MA_11_Social_Policy.tex, frame 3]
\begin{frame}{Average vs. Marginal Tax Rates}
    \showmytikz{../shared/tikz/tikz-marginaltaxrate-lowearners}
    \mycomment{A stylised graph with illustrative purpose. Low income / transfer receivers face prohibitive marginal burdens.}
\end{frame}

%%% [moved from MA_11_Social_Policy.tex, frame 4]
    


\begin{frame}{Taxes on Capital and on Capital Gains}
    
\end{frame}

%%% [moved from MA_11_Social_Policy.tex, frame 5]
\begin{frame}{Welfare Trap}
    \showmytikz{../shared/tikz/tikz-welfaretrap}
\end{frame}

%%% [moved from MA_11_Social_Policy.tex, frame 6]
%%%


\begin{frame}{Transfer Reform}
    \mycomment{
        The abscissa shows income, the ordinate income after transfer. 

        There are some thresholds where income is insufficient. There is a subsistence minium, below which there is starvation. There is also a poverty line, typically set at below 60\% of \textit{median} income.

        Note that this diagram \enquote{zooms into} the lower income echelon.
    }
    \showmytikz[l]{../shared/tikz/tikz-negativetax}
\end{frame}

%%% [moved from MA_11_Social_Policy.tex, frame 7]
\begin{frame}{Transfer Reform}
    \mycomment{
        Under the old policy, everyone whose income is below that threshold receives a benefit, after which there is some heavy taxation to finance that transfer.

        The benefit is targeted, and conditional. The conditions must also be tested and controlled.

        This lifts people typically out of starvation, but keeps them in poverty. 
        The problem is the poverty trap at $P$. There is also the problematic point $T$ after which taxation begins. 
        
        Utility increases strictly to the north with growing disposable income, but slightly to the left because of the disutility of work.
    }
    \showmytikz[l,u,1]{../shared/tikz/tikz-negativetax}
\end{frame}

%%% [moved from MA_11_Social_Policy.tex, frame 8]
\begin{frame}{Transfer Reform}
    \mycomment{
        Those above the threshold must pay enough taxes (transfers) to finance the benefit paid out. Those who receive the benefit do not pay in.

        Note that this diagram does not show the number of benefactors and benefiters. The income distribution is not modelled. The green and red areas are illustrative only. The size of the green area depends on the income distribution.

    }
    \showmytikz[u,l,1,2]{../shared/tikz/tikz-negativetax}
\end{frame}

%%% [moved from MA_11_Social_Policy.tex, frame 9]
\begin{frame}{Transfer Reform}   
    \mycomment{
        The idea is to have a constant marginal tax rate. This avoids the poverty trap. More people will work, and hence contribute. 
        
        This would also allow for a lower marginal tax rate, which is shown here by the slope.

        The point $X$ is a break-even point. Everyone to the left of this will be a net receiver; everyone to the right, a net contributor.
    }
    \showmytikz[u,1,3]{../shared/tikz/tikz-negativetax}
\end{frame}

%%% [moved from MA_11_Social_Policy.tex, frame 10]
\begin{frame}{Transfer Reform}
    \mycomment{
        One way to achieve this is Universal Basic Income (UBI). Everyone receives UBI tax free. Everyone's additional is taxed at the same marginal rate.

        This avoids the poverty trap. At any point, any extra money earned has some utility.

        Typically, UBI is above the subsistence level, but below the poverty line. The exact height is, of course, debatable.
        
    }
    \showmytikz[l,3,4]{../shared/tikz/tikz-negativetax}
\end{frame}

%%% [moved from MA_11_Social_Policy.tex, frame 11]
\begin{frame}{Transfer Reform}
    \mycomment{
        Everyone left of the threshold $X$ is still a net receiver, even though paying tax on earned income.

        You may conclude that for this to work, the red rectangle (UBI times population) should be equal to or smaller the green area. This is true in principle. But we are not modelling the income distribution here. A simple comparison of these areas assumes equal distribution of income. 
    }
    \showmytikz[l,3,4,5]{../shared/tikz/tikz-negativetax}
\end{frame}

%%% [moved from MA_11_Social_Policy.tex, frame 12]
\begin{frame}{Transfer Reform}
    \mycomment{
        Another way to achieve the exact same result is negative income tax (NIT; Friedman), which as a concept is older. 
        
        Everyone whose income is below the break-een point receives payment; everyone above pays in. This break-even point need not be identical with the poverty line; it will depend on the tax rate only, and the tax rate will depend on the actual distribution of income after the policy is implemented.

        NIT is equal to the net effect of UBI.
    }
    \showmytikz[6,7]{../shared/tikz/tikz-negativetax}
\end{frame}

%%% [moved from MA_11_Social_Policy.tex, frame 13]
 

\begin{frame}{Transfer Reform}
    \mycomment{
        Negative Taxes and Universal Basic Income are two different ways to achieve the exact same result. They are equivalent; the only difference is the way of calculation, and how they are institutionally administered (UBI pays upfront).

        The concept of UBI is somehow easier to explain, since no negative numbers appear.
    }
    \showmytikz[l,3,4,5,6,7]{../shared/tikz/tikz-negativetax}
\end{frame}

%%% [moved from MA_11_Social_Policy.tex, frame 14]
\begin{frame}{Transfer Reform}
    \mytext{
    A constant marginal tax - in the form of UBI or Negative Tax - avoids the Poverty Trap.

    It also is administratively vastly more simple. It could at once a whole set of discretionary transfers.

    However, changing an existing complex system is hard. Those who have in the past paid into social security systems cannot easily be expropriated.
    }
\end{frame}

%%% [moved from MA_11_Social_Policy.tex, frame 19]
\begin{frame}{Pensions}
    \showmytikz{../shared/tikz/tikz-pensions.tex}
    \mycomment{
    Pensions depend on
    \begin{enumerate}
        \item \textbf{pension height}
        \item \textbf{payment height}
        \item \textbf{retirement age}
        \item demographic development
        \item life expectancy
    \end{enumerate}
    4) and 5) are exogenous. We are left with three variables: pension height, retirement age, and payment height. Only two of these can be set, determining the third (cf. Tinbergen's Law).
    }
\end{frame}


\begin{frame}{Public or Private Insurance}
    \mycomment{
    A public insurance system has everyone pay a proportion of their income, rather than the real costs of the insurance.
    }
    \showmytikz{../shared/tikz/tikz-publicprivateinsurance}
\end{frame}

\begin{frame}{Public or Private Insurance}
    \mycomment{
        For this to work without state subsidy, high income earners must pay more than their expected payout, subsidising those who pay less.
    }
    \showmytikz[1]{../shared/tikz/tikz-publicprivateinsurance}
\end{frame}

\begin{frame}{Public or Private Insurance}
    \mycomment{
    At a certain income level, people have an incentive to opt out and go for private insurance instead.
    }
    \showmytikz[1,2]{../shared/tikz/tikz-publicprivateinsurance}
\end{frame}

\begin{frame}{Public or Private Insurance}
    \mycomment{
    The result is an efficiency loss, and the public insurance will run a deficit.
    }
    \showmytikz[1,2,3]{../shared/tikz/tikz-publicprivateinsurance}
\end{frame}

\begin{frame}{Public or Private Insurance}
    \mycomment{
    One idea is to add a \enquote{solidarity contribution}, effectively not allowing people to leave the system entirely.

    The other idea would be a state subsidy of the public insurance system.
    }
    \showmytikz[1,2,3,4]{../shared/tikz/tikz-publicprivateinsurance}
\end{frame}


%%% Hochman and Rodgers